\chapter{Ringraziamenti}
{\raggedleft \normalsize \sl "Dimostro gratitudine non perché un altro, \par stimolato dal mio esempio,  mi aiuti più volentieri, \par ma per compiere un'azione bellissima  che mi riempie di gioia; \par sono grato non perché mi conviene, ma perché mi piace."\par 
\bigskip
\--- Seneca, Lettere a Lucilio\par
\bigskip}

\noindent Ogni storia ha un eroe o un'eroina, ma difficilmente il o la protagonista potrebbe compiere la sua missione senza una squadra al suo fianco. Ulisse non avrebbe mai fatto ritorno ad Itaca senza la propria ciurma. Katniss non sarebbe mai uscita viva dall'Arena senza Peeta o Haymitch.

Io non faccio eccezione. Se sono riuscita ad arrivare fino a qui e a compiere il primo piccolo passo verso il coronamento di un sogno, il merito è anche delle persone meravigliose che hanno camminato al mio fianco.
\medskip

Partendo dalle conoscenze più recenti, vorrei ringraziare la mia relatrice, la professoressa Giulia Despali, per la cura, la gentilezza e la costanza con cui mi ha aiutato a redigere questa tesi. Il tempo che mi ha dedicato è stato prezioso: grazie per avermi guidato con pazienza nell’approfondimento non di uno, ma di ben due argomenti che da sempre mi affascinano. 

Vorrei approfittarne per ringraziare anche la professoressa Barbara Lanzoni, che ha permesso il nostro incontro e mi ha fornito le basi di astrofisica. Senza il suo lavoro, la partenza sarebbe stata necessariamente molto più difficoltosa.

Tra le altre insegnanti incontrate nel tempo, un pensiero speciale va alla professoressa Danila Berti, per avermi ricordato il valore della gentilezza in un periodo in cui il buio aveva preso il sopravvento e per la forza d’animo che ha sempre dimostrato. Mi ha insegnato come combattere senza mai arrendersi, nonostante alla fine il suo nemico abbia avuto la meglio.  

Un angolo del mio cuore è dedicato poi a tutta la famiglia del Liceo Scientifico “Augusto Righi” di Cesena, che mi accoglie ancora a braccia aperte ogni volta che torno, per aver reso quelle mura speciali. 

Tra loro, un ringraziamento particolare al professor Francesco Benvenuti, per aver creduto in me anche quando mi sentivo persa e per avermi dato la possibilità di diventare una scienziata più competente, dimostrandosi un insegnante esigente ma non intransigente. Con la sua passione ha acceso la mia, e La ricorderò sempre con affetto. 

Tra i docenti universitari vorrei rivolgere un ringraziamento a quanti mi hanno dimostrato che, tra chi spiega in cattedra e chi siede tra i banchi, non deve per forza esserci una distanza abissale o un rapporto impersonale. Un grazie speciale alla professoressa Sara Valentinetti, che più di una volta mi ha aperto le porte del suo laboratorio, accogliendomi con tutti i miei dubbi e le mie fragilità, e condividendo le sue. 

Grazie anche a tutti i tutor e i tecnici di laboratorio, per la pazienza, l'aiuto nel sostituire i diodi bruciati e per le preziose dispense.
\medskip

Tra chi ha condiviso i banchi con me, in primo luogo vorrei ringraziare tutti i membri del gruppo \emph{eh?! nononono, basta statistica}. Mi avete dimostrato, con ogni pranzo e ogni serata, che non servono grandi cose per essere felici. Grazie per aver reso speciali i piccoli gesti quotidiani e per l’aiuto reciproco che ci siamo dati, nello studio come nelle piccole difficoltà di ogni giorno. 

In particolare, grazie a Bianca, che per tre anni è stata la mia fedele compagna nelle lunghe ore in laboratorio e nella stesura delle relazioni. Averti al mio fianco me le ha fatte pesare un po’ meno. Anche se negli ultimi anni non ci siamo viste ogni giorno, hai sempre trovato il modo di farti sentire vicina quando ne avevo più bisogno. 

A Giulia, grazie per le chiacchiere sul passato e sul futuro e per aver messo a disposizione il tuo frigorifero quando il mio si è rotto.

A Riccardo, grazie per tutto il supporto e il sostegno. Dall’esame di  Algebra Lineare a quello di Meccanica Quantistica, in questi tre anni abbiamo fatto tutto insieme. Avrei dovuto capirlo quando abbiamo iniziato a chiacchierare di Seneca e Montale che saresti diventato uno dei miei migliori amici. Grazie per essere stato al mio fianco, sempre, e per avermi spronato a \emph{nun mullà!} quando ho vacillato. 

Una parte importante in questa esperienza di crescita l’hanno avuta anche le mie coinquiline: grazie per avermi dato un luogo sicuro in cui tornare nella grande metropoli di Bologna, per tutte le serate sul nostro divano, che ormai ne sa una più del diavolo, per i post-it colorati lasciati sul tavolo e per le cioccolate calde condivise in inverno. 

In particolare, grazie a Marianna, che ha saputo vedere oltre il finto sorriso indossato in un momento di inquietudine profonda. Mi hai aiutata ogni giorno, standomi accanto senza invadenza, a volermi bene e a raccogliere il coraggio necessario ad agire di conseguenza. 
\medskip

Se i miei amici e le mie amiche sono stati come una seconda famiglia a cento chilometri da casa, il ringraziamento più sentito va ai miei genitori. Mi avete insegnato a lottare per ciò in cui credo e ad amare la cultura, nel senso più ampio del termine. Grazie, per avermi sostenuta incondizionatamente nella mia scelta di laurearmi in Fisica, anche se talvolta non l’avete compresa. 

Grazie a mamma, perché non mi ha mai nascosto come essere una donna nella nostra società non sia facile, ma nonostante questo mi ha sempre spronato a lavorare sodo per qualsiasi obiettivo mi fossi posta, indipendentemente dalle convenzioni sociali. Grazie per tutti i consigli di lettura, cucina, moda e vita in generale. E grazie, per avermi insegnato il valore dell’onestà, qualunque cosa accada. 

Grazie a papà, che mi ha sempre accompagnato ovunque volessi andare, indipendentemente dalla distanza. E grazie per la passione condivisa per tutto il mondo scientifico. 

Vorrei ringraziare anche mia sorella, perché, nonostante su molte cose siamo agli antipodi, riusciamo comunque a trovare delle passioni condivise e a volerci bene a modo nostro. 

Grazie ad Ariel, che ha dormito sui miei appunti durante le lunghe ore di studio a casa, tenendomi compagnia con costanza. Con i poteri conferitomi (ovvero neanche uno) ti dichiaro Laureata in Fisica ad honorem. 

Grazie a tutti i parenti, ed in particolare a nonna, per aver provato a capire cosa faccio e avermi sostenuto anche quando non ci sono riusciti. 

Grazie a tutti gli amici storici, ormai annoverati in famiglia, perché nonostante la distanza che separa alcuni di noi, non mi avete mai fatto sentire sola.

E grazie a tutti coloro che hanno condiviso con me una parte del viaggio, anche se non ne hanno visto il finale, per esserci stati, almeno in quel tratto. 
\medskip

Vorrei infine ringraziare sentitamente tutte le donne, e in particolar modo quelle di scienza, che nel corso della Storia si sono battute con determinazione e costanza per i nostri diritti e per il riconoscimento dei nostri meriti. Non sarei qui oggi senza di voi.
\medskip 

Infine, grazie a te lettore o lettrice: spero tu abbia imparato qualcosa in più rispetto a quando ti sei imbattuta in questo elaborato, che sia scientifica o meno. Grazie, per essere rimasta con me, fin proprio alla fine.

